% !TeX root = ../analisi_2.tex

\section{Integrazione multipla}

\subsection{Domini normali}

\begin{oss}
    Data una funzione $f : A \sube \R^2 \rarr \R$, il suo grafico è la superficie
    \[
        z = f(x,y), \qquad (x,y) \in A.
    \]
    Se $f \ge 0$, la regione di spazio compresa tra $A$ e il grafico di $f$ si chiama \textbf{cilindroide} di $f$:
    \[
        C_f = \{ (x,y,z) \in \R^3 : (x,y) \in A,\; 0 \le z \le f(x,y) \}.
    \]
\end{oss}

\pencilnote{C'è il problema delle ``partizioni'': sono ``curve'' (non rettangoli).
    Come si risolve?}

La risposta sta nel lavorare con insiemi del piano di un tipo particolare, che possiamo ``affettare'' in modo semplice: i \textbf{domini normali}.

\begin{definition}[Dominio normale rispetto a $x$]
    Si dice \textbf{dominio normale rispetto a $x$} (in $\R^2$) l'insieme
    \[
        D = \{ (x,y) \in \R^2 : a \le x \le b,\; \alpha(x) \le y \le \beta(x) \}
    \]
    con $\alpha \le \beta$ e $\alpha$, $\beta$ continue.
\end{definition}

\pencilnote{Se ``taglio'' in verticale ho ``un pezzo'', se ``taglio'' in orizzontale potrei averne di più.}

\begin{definition}[Dominio normale rispetto a $y$]
    Un insieme si dice \textbf{normale rispetto a $y$} se
    \[
        D = \{ (x,y) \in \R^2 : c \le y \le d,\; \gamma(y) \le x \le \delta(y) \}
    \]
    con $\gamma \le \delta$ e $\gamma$, $\delta$ continue.
\end{definition}

\begin{definition}[Area di un dominio normale]
    L'\textbf{area} di $D$ (ovvero la sua \emph{misura}) si indica con $m(D)$ ed è data da
    \[
        m(D) = \int_a^b \big( \beta(x) - \alpha(x) \big) \,dx
        \qquad \text{se $D$ è normale rispetto a } x,
    \]
    \[
        m(D) = \int_c^d \big( \delta(y) - \gamma(y) \big) \,dy
        \qquad \text{se $D$ è normale rispetto a } y.
    \]
\end{definition}

\pencilnote{L'integranda è positiva poiché $\beta \ge \alpha$ (e $\delta \ge \gamma$).}

\subsection{Partizioni in domini normali}

\begin{definition}[Partizione in domini normali]
    Sia $D$ un dominio normale. Si chiama \textbf{partizione di $D$ in domini normali} la seguente famiglia
    \[
        \mathcal{P} = \{ D_1, D_2, \dots, D_N \}
    \]
    tale che
    \begin{enumerate}[1)]
        \item $D_i$ sono tutti normali $\forall i$;
        \item $\displaystyle D = \bigcup_{i=1}^{N} D_i$;
        \item $\overset{\circ}{D_i} \cap \overset{\circ}{D_j} = \emptyset \quad \forall i \neq j$.
    \end{enumerate}
\end{definition}

\begin{lemma*}[senza dimostrazione]
    Sia $D$ un dominio normale. Allora esiste una partizione di $D$ in domini normali (fitta quanto si vuole).
\end{lemma*}

%   \include{Novembre/16-11}
%   \include{Novembre/19-11}

\subsection{Integrali doppi}

Sia $f : D \sube \R^2 \rarr \R$, con $D$ un dominio normale del piano, e sia $f$ \textbf{limitata}.
Per il lemma visto in precedenza esiste una partizione di $D$ in domini normali
\[
    \mathcal{P} = \{ D_1, D_2, \dots, D_N \}.
\]

Poiché $f$ è limitata, possiamo porre
\[
    m_i = \inf_{D_i} f(x,y) < +\infty, \qquad
    M_i = \sup_{D_i} f(x,y) < +\infty .
\]

Allora possiamo definire le seguenti \textbf{somme integrali}:
\begin{align*}
    s(\mathcal{P}, f) &= \sum_{i=1}^{N} m(D_i) \cdot m_i
    && \text{(somma integrale inferiore relativa a $f$ rispetto a $\mathcal{P}$)} \\
    S(\mathcal{P}, f) &= \sum_{i=1}^{N} m(D_i) \cdot M_i
    && \text{(somma integrale superiore relativa a $f$ rispetto a $\mathcal{P}$)}
\end{align*}
dove $m(D_i)$ indica l'area di $D_i$.

\begin{lemma}
    Siano $\mathcal{P}$ e $\mathcal{Q}$ due partizioni di $D$ in domini normali. Allora
    \[
        s(\mathcal{P}, f) \le S(\mathcal{Q}, f).
    \]
\end{lemma}

Siano
\begin{align*}
    A &= \{ s(\mathcal{P}, f) \,,\; \mathcal{P} \text{ partizione} \} \sube \R, \\
    B &= \{ S(\mathcal{P}, f) \,,\; \mathcal{P} \text{ partizione} \} \sube \R.
\end{align*}
Per il lemma, $A$ e $B$ sono \textbf{separati}.

\begin{definition}[Integrale doppio]
    Se $A$ e $B$ sono \textbf{contigui}, ovvero $\sup A = \inf B = c$, allora l'elemento di separazione $c$ prende il nome di
    \textbf{integrale doppio} di $f$ su $D$ e si indica con
    \[
        \iint_D f(x,y) \,dx\,dy .
    \]
    In tal caso $f$ si dice \textbf{integrabile secondo Riemann} su $D$.
\end{definition}

\begin{oss}[Misura di Peano--Jordan]
    La misura di Riemann è la misura di Peano--Jordan.

    \pencilnote{Divido in due ``gabbie'': una esterna e una interna.
        La gabbia interna è fatta di rettangoli e, detta $A(R_{\mathcal{P}}) = \sum_i A(R_i)$ l'area di un plurirettangolo contenuto in $D$, si pone
        \[
            m_i = \sup \{ A(R_{\mathcal{P}}) \,;\; R_{\mathcal{P}} \sube D \} \qquad \text{(misura interna)}
        \]
        Analogamente per l'esterna:
        \[
            m_e = \inf \{ A(R_{\mathcal{P}}) \,;\; R_{\mathcal{P}} \supseteq D \} \qquad \text{(misura esterna)}
        \]
        Se $m_i = m_e$, il separatore è detto misura di Peano--Jordan.
        Lebesgue invece ``approssima'' con compatti interni e aperti esterni.}
\end{oss}

\subsection{Proprietà dell'integrale doppio}

\begin{enumerate}[1)]
    \item \textbf{Linearità:}
          \[
              \iint_D (\alpha f + \beta g) \,dx\,dy = \alpha \iint_D f \,dx\,dy + \beta \iint_D g \,dx\,dy .
          \]
    \item \textbf{Monotonia:} se $f \le g$ allora
          \[
              \iint_D f \,dx\,dy \le \iint_D g \,dx\,dy .
          \]
    \item \textbf{Additività:} se $D = \displaystyle\bigcup_{i=1}^{k} D_i$ con $D_i$ normale $\forall i$ e $\overset{\circ}{D_i} \cap \overset{\circ}{D_j} = \emptyset$ per $i \neq j$, allora
          \[
              \iint_D f \,dx\,dy = \sum_{i=1}^{k} \iint_{D_i} f \,dx\,dy .
          \]
\end{enumerate}

\begin{theorem*}[Integrabilità delle funzioni continue]
    Se $f$ è continua, allora è integrabile.
\end{theorem*}

\begin{prop*}[Teorema della media (stima)]
    Se $f$ è continua, allora
    \[
        m(D) \cdot \min_D f \;\le\; \iint_D f \,dx\,dy \;\le\; m(D) \cdot \max_D f .
    \]
\end{prop*}

\subsection{Formule di riduzione}

\begin{theorem*}[Formule di riduzione (senza dimostrazione)]
    Sia $D$ normale e sia $f \colon D \sube \R^2 \rarr \R$ continua.

    Se $D$ è normale rispetto a $x$, ovvero $D = \{ a \le x \le b,\ \alpha(x) \le y \le \beta(x) \}$, risulta
    \[
        \iint_D f(x,y) \,dx\,dy = \int_a^b \left[ \int_{\alpha(x)}^{\beta(x)} f(x,y) \,dy \right] dx .
    \]

    Se invece $D$ è normale rispetto a $y$, ovvero $D = \{ c \le y \le d,\ \gamma(y) \le x \le \delta(y) \}$, si ha
    \[
        \iint_D f(x,y) \,dx\,dy = \int_c^d \left[ \int_{\gamma(y)}^{\delta(y)} f(x,y) \,dx \right] dy .
    \]
\end{theorem*}

\pencilnote{Fissata una $x$, la zona colorata è un rettangoloide.}

\begin{esercizio}
    Calcolare
    \[
        \iint_E x e^y \,dx\,dy, \qquad E = \{ 0 \le x \le 1,\ x^2 + y^2 \le 1,\ y \ge 0 \}.
    \]
\end{esercizio}

\begin{solution}
    $E$ è il quarto di cerchio unitario nel primo quadrante, quindi può essere scritto in entrambi i modi:
    \[
        E = \{ x \in [0,1],\ 0 \le y \le \sqrt{1 - x^2} \}
        \qquad \text{oppure} \qquad
        E = \{ y \in [0,1],\ 0 \le x \le \sqrt{1 - y^2} \}.
    \]

    \medskip
    \textbf{Prima parametrizzazione} (normale rispetto a $x$):
    \[
        \iint_E x e^y \,dx\,dy
        = \int_0^1 \left[ x \int_0^{\sqrt{1-x^2}} e^y \,dy \right] dx
        = \int_0^1 x \left[ e^{\sqrt{1-x^2}} - e^0 \right] dx
        = \int_0^1 x\, e^{\sqrt{1-x^2}} \,dx - \int_0^1 x \,dx .
    \]
    Nel primo integrale poniamo $t = \sqrt{1-x^2}$, da cui
    \[
        dt = \frac{-2x}{2\sqrt{1-x^2}} \,dx = -\frac{x}{t} \,dx
        \quad \Longrightarrow \quad -t \,dt = x \,dx,
    \]
    con $x = 0 \mapsto t = 1$ e $x = 1 \mapsto t = 0$. Quindi
    \[
        \int_0^1 x\, e^{\sqrt{1-x^2}} \,dx = \int_1^0 e^t (-t) \,dt = \int_0^1 t e^t \,dt
        = \Big[ e^t (t - 1) \Big]_0^1 = 1,
    \]
    e dunque
    \[
        \iint_E x e^y \,dx\,dy = 1 - \frac{1}{2} = \frac{1}{2}.
    \]

    \medskip
    \textbf{Seconda parametrizzazione} (``l'altra'', normale rispetto a $y$):
    \[
        \iint_E x e^y \,dx\,dy
        = \int_0^1 e^y \left[ \int_0^{\sqrt{1-y^2}} x \,dx \right] dy
        = \int_0^1 e^y \left[ \frac{x^2}{2} \right]_0^{\sqrt{1-y^2}} dy
        = \int_0^1 e^y \, \frac{1 - y^2}{2} \,dy .
    \]
    Integrando due volte per parti:
    \[
        \frac{1}{2} \int_0^1 e^y (1 - y^2) \,dy
        = \frac{1}{2} \Big[ e^y (1 - y^2) + 2 e^y (y - 1) \Big]_0^1
        = \frac{1}{2} \Big[ e^y (-y^2 + 2y - 1) \Big]_0^1
        = \frac{1}{2} \big( 0 + 1 \big) = \frac{1}{2},
    \]
    in accordo con il primo calcolo.
\end{solution}

\begin{esercizio}
    Calcolare
    \[
        \iint_D \frac{\sin y}{y} \,dx\,dy, \qquad D = \{ 0 \le x \le 1,\ \pi x \le y \le \pi \}.
    \]
\end{esercizio}

\begin{solution}
    Integrando prima in $y$ si avrebbe
    \[
        \int_0^1 \left[ \int_{\pi x}^{\pi} \frac{\sin y}{y} \,dy \right] dx = \; ??
    \]
    ma $\dfrac{\sin y}{y}$ non ha una primitiva elementare. Conviene allora descrivere $D$ come dominio normale rispetto a $y$:
    \[
        D = \left\{ 0 \le y \le \pi,\ 0 \le x \le \frac{y}{\pi} \right\}.
    \]
    Risulta
    \[
        \iint_D \frac{\sin y}{y} \,dx\,dy
        = \int_0^{\pi} \left[ \frac{\sin y}{y} \int_0^{y/\pi} dx \right] dy
        = \int_0^{\pi} \frac{\sin y}{\pi} \,dy
        = \left[ -\frac{\cos y}{\pi} \right]_0^{\pi}
        = \frac{2}{\pi} .
    \]
\end{solution}
